Die Cola hat etwa die Dichte und die spezifische Wärme von Wasser, ihre Masse ist $m_\text{C} = \rho\,V = \mCo$. Wärmeverluste sind vernachlässigt. Das Eis erwärmt sich auf $\vartheta_0 = \TnO$ ($c_\text{E} = \cEO$), schmilzt ($L_\text{f} = \LfEO$, FoTaBe S.~190) und das Schmelzwasser erwärmt sich auf die Endtemperatur $\vartheta$; die Cola kühlt von $\vartheta_\text{C}$ auf $\vartheta$ ab:
\begin{align}
 m_\text{E}\left(c_\text{E}\left(\vartheta_0-\vartheta_\text{E}\right)+L_\text{f}+c_\text{W}\left(\vartheta-\vartheta_0\right)\right) &= m_\text{C}c_\text{W}\left(\vartheta_\text{C}-\vartheta\right)
\end{align}
Aufgelöst nach $\vartheta$:
\begin{align}
 \vartheta &= \TmF \\
   &= \frac{\mCo\cdot\cW\cdot\TCO+\mE\cdot\cW\cdot\TnO-\mE\cdot\left(\cE\cdot\left(\TnO-\left(\TEO\right)\right)+\LfE\right)}{\left(\mCo+\mE\right)\cdot\cW} \\
   &= \Tm \approx \result{\TmP{0}{2}}
\end{align}
